%%%%%%%%%%%%%%%%%%%%%%%%%%%%%%%%%%%%%%%%%%%%%%%%%%%%%%%%%%%%%%%%%%%%%%
% Lớp 11 | Dãy số. Cấp số cộng. Cấp số nhân | Cấp số cộng
% Dạng 1D2?2-2: Nhận diện cấp số cộng, công sai 𝑑
% Số câu: 88  (sinh tự động bởi tach_cau.py; từ nay thầy sửa trực tiếp file này)
%%%%%%%%%%%%%%%%%%%%%%%%%%%%%%%%%%%%%%%%%%%%%%%%%%%%%%%%%%%%%%%%%%%%%%

%%%%%%%%%%%%%%%% TRẮC NGHIỆM 4 PHƯƠNG ÁN %%%%%%%%%%%%%%%%

% nguồn: Ok_Toan11_TracNghiem_Lan2
\begin{ex}%[1D2N2-2]%[ma:c4fd3436]
	Cho cấp số  cộng $1$; $1$; $1$; $1$; $\ldots$. Công sai của cấp số cộng trên là
	 \choice
		 {\True $d=0$}
		 {$d=1$}
		{$d=-1$} 
		 {$d\in\varnothing $}
	 \loigiai{
	 Cấp số cộng có công sai $d=1-1=0$.}
\end{ex}

% nguồn: Ok_Toan11_TracNghiem_Lan2
\begin{ex}%[1D2N2-2]%[ma:4673f88f]
	Trong các dãy số $(u_n)$ được cho dưới đây, dãy số  nào là cấp số cộng?
	 \choice
		{\True $u_n=7-3n$} 
		{$u_n=7-{3^n}$} 
		{$u_n=\dfrac{7}{3n}$} 
		{$u_n={7\cdot{3}^n}$} 
	\loigiai{ 
		Dãy $(u_n)$ là cấp số cộng $\Leftrightarrow u_n=an+b$ ($a, b$ là  hằng số).
	} 
\end{ex}

% nguồn: Ok_Toan11_TracNghiem_Lan2
\begin{ex}%[1D2N2-2]%[ma:147a5889]
	 Cho cấp số cộng có số hạng đầu $u_3=-5$, $u_5=3$. Công sai của cấp số cộng đã cho là
	\choice 
		{$d=2$} 
		 {\True $d=4$}
		{$d=8$} 
		{$d=1$} 
	\loigiai{ 
	Theo định nghĩa cấp  số cộng $u_5-u_3=2d\Leftrightarrow d=4$.}
\end{ex}

% nguồn: Ok_Toan11_TracNghiem_Lan2
\begin{ex}%[1D2N2-2]%[ma:32ed3af0]
	 Trong các dãy số có công thức tổng quát sau, dãy số nào là cấp số cộng?
	 \choice
		 {$u_n=2021^n$}
		 {\True $u_n=2n+2021$}
		{$u_n=\dfrac{2}{n+2021}$} 
		 {$u_n=n^2-2$}
	\loigiai{ 
	Với mỗi $n \in \mathbb{N}^*$ bất kì, $u_n=2n+2021$ thì $u_{n+1}=2(n+1)+2021=u_n+2$, như  vậy dãy số này là một cấp số cộng.}
\end{ex}

% nguồn: Ok_Toan11_TracNghiem_Lan2
\begin{ex}%[1D2N2-2]%[ma:1d72d8c7]
	Trong các dãy  số sau, dãy nào là cấp số cộng?
	\choice 
		{$1$; $2$; $4$; $7$; $11$;  $\ldots$}
		{\True  $\dfrac{1}{3}$; $\dfrac{2}{3}$; $1$; $\dfrac{4}{3}$; $\dfrac{5}{3}$; $\ldots$}
		{$-5$; $-3$; $-1$;  $3$; $5$; $\ldots$}
		{$1$; $2$; $4$; $8$;  $16$; $\ldots$}
	 \loigiai{
	Dãy số $\dfrac{1}{3}$; $\dfrac{2}{3}$; $1$; $\dfrac{4}{3}$; $\dfrac{5}{3}$; $\ldots$ thỏa mãn số hạng đứng trước cộng thêm $\dfrac{1}{3}$ thì được số hạng tiếp theo nên là  cấp số cộng với công sai $d=\dfrac{1}{3}$.}
\end{ex}

% nguồn: Ok_Toan11_TracNghiem_Lan2
\begin{ex}%[1D2N2-2]%[ma:ef95a38c]
	Cho cấp số cộng $\left(u_n\right)$ có số  hạng tổng quát $u_n=2n+5$, $\forall n \in \mathbb{N}^*$. Công sai $d$ của cấp số cộng bằng
	\choice 
		{$5$} 
		 {\True $2$}
		 {$-2$}
		{$-5$} 
	 \loigiai{
	Với mỗi $n \in \mathbb{N}^*$  bất kì, ta có $d=u_{n+1}-u_n=2\left(n+1\right)+5-\left(2n+5\right)=2$.}
\end{ex}

% nguồn: Ok_Toan11_TracNghiem_Lan2
\begin{ex}%[1D2N2-2]%[ma:2bbabee2]
	Cho cấp số cộng  $\left(u_n\right)$ với $u_1=2$ và $u_2=6$. Công sai của cấp số cộng đã cho bằng
	\choice 
		{\True  $4$}
		 {$-4$}
		 {$8$}
		{$3$} 
	\loigiai{ 
	 Ta có $u_2=6\Leftrightarrow 6=u_1+d \Leftrightarrow d=4$.}
\end{ex}

% nguồn: Ok_Toan11_TracNghiem_Lan2
\begin{ex}%[1D2N2-2]%[ma:ca60e066]
	Cho dãy $(u_n)$ là một cấp số cộng với số hạng đầu $2$ và số hạng thứ $36$  là $72$. Công sai của cấp số cộng $(u_n)$ là
	 \choice
		 {$3$}
		 {$-2$}
		 {\True $2$}
		{$\dfrac{1}{2}$} 
	 \loigiai{
		Ta có $u_{36}  = u_1+35\cdot d$ mà $u_{36}=72$, $u_1=2$ nên ta có $72 = 2+35\cdot d \Leftrightarrow d=2$.
	} 
\end{ex}

% nguồn: Ok_Toan11_TracNghiem_Lan2
\begin{ex}%[1D2N2-2]%[ma:11e3f693]
	Cho cấp số cộng $\left(u_n\right)$ với $u_n=5-2n$, $\forall n \in \mathbb{N}^*$. Tìm công sai của  cấp số cộng
	 \choice
		 {$d=3$}
		{\True  $d=-2$}
		{$d=1$} 
		 {$d=2$}
	 \loigiai{
	Với mỗi $n \in \mathbb{N}^*$ bất kì, ta có  \[u_{n+1}-u_n=\left(5-2\left(n+1\right)\right)-\left(5-2n\right)=5-2n-2-5+2n=-2\Rightarrow d=-2.\]}
\end{ex}

% nguồn: Ok_Toan11_TracNghiem_Lan2
\begin{ex}%[1D2N2-2]%[ma:902f3cf9]
	Trong các dãy số sau, dãy số nào là một cấp  số cộng?
	\choice 
		{$1$; $-2$; $-4$; $-6$; $-8$} 
		{$1$; $-3$; $-6$; $-9$;  $-12$}
		{\True  $1$; $-3$; $-7$; $-11$; $-15$}
		{$1$; $-3$; $-5$;  $-7$; $-9$}
	 \loigiai{
		 Dãy số $\left(u_n\right)$ có tính chất $u_{n+1}=u_n+d$ thì được gọi là một cấp số cộng.\\
	Ta thấy dãy số $1$; $-3$; $-7$; $-11$; $-15$ là một cấp số cộng có  số hạng đầu là $1$ và công sai bằng $-4$.}
\end{ex}

% nguồn: Ok_Toan11_TracNghiem_Lan2
\begin{ex}%[1D2N2-2]%[ma:99256bfd]
	Trong  các dãy số $(u_n)$ sau, dãy số nào là một cấp số cộng?
	 \choice
		{$1;  \,-2; \,-4; \,-6; \,-8$}
		{$1; \,-3; \,-6; \,-9;  \,-12$}
		{\True $1; \,-3; \,-7; \,-11; \,-15$} 
		 {$1; \,-3; \,-5; \,-7; \,-9$}
	\loigiai{ 
		Dãy số $(u_n)$ là cấp số cộng khi $u_{n+1}=u_n+d$  với $d$ không đổi. \\
		Ta thấy dãy $1; \,-3; \,-7; \,-11; \,-15$ có công sai $d=-4$ nên là  cấp số cộng.
	 }
\end{ex}

% nguồn: Ok_Toan11_TracNghiem_Lan2
\begin{ex}%[1D2N2-2]%[ma:9c4f5d3a]
	Cho cấp số cộng $\left(u_n\right)$ với $u_n=5-2n$, $\forall  n \in \mathbb{N}^*$. Tìm công sai của cấp số cộng.
	 \choice
		 {$d=3$}
		{\True  $d=-2$}
		 {$d=1$}
		{$d=2$} 
	 \loigiai{
		Với mỗi $n \in \mathbb{N}^*$ bất kì, ta có $u_n=5-2n\Rightarrow{u_{n+1}}=5-2\left(n+1\right)=3-2n$.\\ 
	 Suy ra $d=u_{n+1}-u_n=3-2n-\left(5-2n\right)=-2$.}
\end{ex}

% nguồn: Ok_Toan11_TracNghiem_Lan2
\begin{ex}%[1D2N2-2]%[ma:08fd1cbf]
	Cho cấp số cộng $(u_n)$ có $u_1=-5$; $u_8=30$. Công sai của cấp số cộng  bằng
	\choice 
		 {$4$}
		{\True $5$} 
		{$6$} 
		{$3$} 
	\loigiai{ 
		Gọi công sai của cấp số cộng  là $d$ khi đó ta có $u_8=u_1+7\cdot d \Leftrightarrow 30=-5+7\cdot d \Leftrightarrow d=5$.
	} 
\end{ex}

% nguồn: Ok_Toan11_TracNghiem_Lan2
\begin{ex}%[1D2N2-2]%[ma:49bdcd0f]
	Cho cấp số cộng $(u_n)$ có $u_2=8$;  $u_5=17$. Công sai $d$ bằng
	\choice 
		 {$-3$}
		 {$-5$}
		{\True  $3$}
		{$5$} 
	 \loigiai{
		Theo giả thiết ta có $u_2=8; u_5=17 \Rightarrow  \heva{&u_1+d=8\\&u_1+4\cdot d=17} \Leftrightarrow \heva{&d=3\\&u_1=5.}$
	} 
\end{ex}

% nguồn: Ok_Toan11_TracNghiem_Lan2
\begin{ex}%[1D2N2-2]%[ma:75721084]
	Cho cấp số cộng $(u_n)$  với $u_1=2$ và $u_2=7$. Công sai của cấp số cộng đã cho bằng
	\choice 
		{\True$5$} 
		{$\dfrac{2}{7}$} 
		 {$-5$}
		{$\dfrac{7}{2}$} 
	 \loigiai{Ta có $u_2-u_1=7-2=5$.
	} 
\end{ex}

% nguồn: Ok_Toan11_TracNghiem_Lan2
\begin{ex}%[1D2N2-2]%[ma:c22c7b2f]
	Trong các dãy số $(u_n)$ được cho dưới đây,  dãy số nào là cấp số cộng?
	 \choice
		{$u_n=(-1)^n(2n+1)$} 
		 {$u_n=\sin \left(\dfrac{\pi}{n}\right)$}
		{\True $\left\{\begin{aligned}
		 u_1&=1\\
		u_n  &=u_{n-1}-1
\end{aligned}\right.$}
{$\left\{\begin{aligned}
 u_1&=1\\
 u_n &=2u_{n-1}
\end{aligned}\right.$}
 \loigiai{
	Một dãy $(u_n)$ là cấp số cộng khi  tồn tại hằng số $d$ sao cho
	\[ 
	u_n=u_{n-1}+d \quad \text{  (với $d$ không đổi và $n \ge 2$)}.
	 \]
	Xét  dãy số $u_n=3n+1$. Ta có\\
	$u_n=u_{n-1}-1$ nghĩa  là công sai $d=-1$, không đổi.\\ Vậy $u_n=3n+1$ là cấp số cộng.
 }
\end{ex}

% nguồn: Ok_Toan11_TracNghiem_Lan2
\begin{ex}%[1D2N2-2]%[ma:bf0e0804]
	Cho cấp số cộng $5$;  $10$; $15$; $20$; $25$; $\ldots$. Khi đó công sai $d$ bằng
	\choice 
		{$d=1$} 
		 {$d=2$}
		 {\True $d=5$}
		 {$d=10$}
	\loigiai{ 
	Cấp số cộng đã cho có công sai  $d=5$.}
\end{ex}

% nguồn: Ok_Toan11_TracNghiem_Lan2
\begin{ex}%[1D2N2-2]%[ma:5a79cdbc]
	Cho cấp số cộng $(u_n)$ với $u_1=2$  và $u_2=6$. Công sai của cấp số cộng đã cho bằng
	\choice 
		{\True$4$} 
		 {$-4$}
		 {$8$}
		 {$3$}
	 \loigiai{Ta có $d=u_2-u_1=6-2=4$.
	} 
\end{ex}

% nguồn: Ok_Toan11_TracNghiem_Lan2
\begin{ex}%[1D2N2-2]%[ma:edcb329c]
	Trong  các dãy số sau, dãy số nào không phải là cấp số cộng?
	 \choice
		{ $1;\, 1;  \, 1; \, 1; \, 1$}
		{ $\dfrac{1}{2}; \, \dfrac{3}{2};  \, \dfrac{5}{2}; \, \dfrac{7}{2}; \, \dfrac{9}{2}$}
		{ $-8; \,-6; \,-4; \,-2;  \, 0$}
		{\True  $3; \, 1; \,-1; \,-2; \,-4$}
	\loigiai{ 
		Xét dãy số $3; \, 1;  \,-1; \,-2; \,-4$.\\
		Ta có $u_2-u_1=-2$, nhưng $u_4-u_3=-1$. 
		không phải cấp số cộng vì hiệu không bằng nhau.\\Nên dãy số này không phải là cấp số  cộng.
	 }
\end{ex}

% nguồn: Ok_Toan11_TracNghiem_Lan2
\begin{ex}%[1D2N2-2]%[ma:4cc14dcd]
	Trong các dãy số  $\left(u_n\right)$ sau đây, dãy số nào là cấp số cộng?
	\choice 
		 {$\heva{&u_1=3\\&{u_{n+1}}=2u_n+1}$}
		{\True  $\heva{& {u_1}=-1\\& {u_{n+1}}-u_n=2}$}
		 {$\heva{&{u_1}=1\\& {u_{n+1}}=u_n^3-1}$}
		 {$\heva{& {u_1}=1\\&{u_{n+1}}=u_n+n}$}
	\loigiai{ 
		Xét $\heva{&u_1=3\\&{u_{n+1}}=2u_n+1}$ là $u_2=7$, $u_3=15$ vì $u_2-u_1\ne u_3-u_2$ do đó $\left(u_n\right)$ không  phải là cấp số cộng.\\
		Xét $\heva{& {u_1}=-1\\& {u_{n+1}}-u_n=2}$ là theo giả thiết ta có do đó $\left(u_n\right)$ là  cấp số cộng.\\
		Xét $\heva{&{u_1}=1\\& {u_{n+1}}=u_n^3-1}$ là $u_2=0$, $u_3=-1$, $u_4=-2$, $u_5=-9$ do đó $\left(u_n\right)$ không phải  là cấp số cộng.\\
	 Xét $\heva{&{u_1}=1\\&{u_{n+1}}=u_n+n}$ là $u_2=2$, $u_3=4$ vì $u_2-u_1\ne\,u_3-u_2$ do đó $\left(u_n\right)$ không phải là cấp số cộng.}
\end{ex}

% nguồn: Ok_Toan11_TracNghiem_Lan2
\begin{ex}%[1D2N2-2]%[ma:2c989cb9]
	Trong các dãy số sau, dãy  số nào là một cấp số cộng?
	\choice 
		{\True $1$; $4$; $7$;  $10$; $13$}
		{$1$;  $3$; $1$; $3$; $1$}
		{$1$; $3$; $7$; $10$; $13$} 
		{$2$, $4$,  $5$, $6$, $8$}
	\loigiai{ 
		Dãy số $\left(u_n\right)$ có tính chất $u_{n+1}=u_n+d$, $\forall n  \in \mathbb{N}^*$ thì được gọi là một cấp số cộng.\\
	Ta thấy dãy số $1$; $4$; $7$; $10$; $13$ là một cấp số cộng có số hạng đầu là $u_1=1$ và công sai bằng  $d=3$.}
\end{ex}

% nguồn: Ok_Toan11_TracNghiem_Lan2
\begin{ex}%[1D2N2-2]%[ma:79f6735b]
	Trong các dãy  số $(u_n)$ được cho dưới đây, dãy số nào không phải là cấp số cộng?
	 \choice
		{$u_n=-4n+9$} 
		{$u_n=-2n+19$} 
		 {$u_n=-2n-21$}
		{\True  $u_n=-{2^n}+15$}
	\loigiai{ 
		Dãy số $u_n=-{2^n}+15$ không  có dạng $a\cdot n+b$ nên có không phải là cấp số cộng.
	 }
\end{ex}

% nguồn: Ok_Toan11_TracNghiem_Lan2
\begin{ex}%[1D2N2-2]%[ma:119d9dfb]
	Cho cấp số cộng  $(u_n)$ với $u_n=5-2 n$. Công sai của cấp số cộng là
	\choice 
		{$d=3$} 
		 {\True $d=-2$}
		{$d=1$} 
		 {$d=2$}
	\loigiai{ 
		Công sai của cấp số cộng  được tính bởi\\
		 \[
		d=u_{n+1}-u_n=\left[5-2(n+1)  \right]-\left(5-2n\right)=5-2 n-2-5+2 n=-2.
		 \]\\
		Vậy công sai  là $d=-2$.
	} 
\end{ex}

% nguồn: Ok_Toan11_TracNghiem_Lan2
\begin{ex}%[1D2N2-2]%[ma:3af53fdd]
	Cho cấp số cộng $(u_n)$ với $u_1=1$  và $u_2=4$. Công sai của cấp số cộng đã cho bằng
	 \choice
		 {$4$}
		{$-3$} 
		{\True  $3$}
		 {$5$}
	 \loigiai{Ta có $d=u_2-u_1=4-1=3$.
	 }
\end{ex}

% nguồn: Ok_Toan11_TracNghiem_Lan2
\begin{ex}%[1D2N2-2]%[ma:0ff6bf73]
	Cho cấp số cộng $\left(u_n\right)$ có $u_1=-3$,  $u_6=27$. Tính công sai $d$ .
	\choice 
		 {$d=7$}
		{$d=5$} 
		 {$d=8$}
		 {\True $d=6$}
	 \loigiai{
	Ta có $u_6=u_1+5d=27\Rightarrow  d=6$.}
\end{ex}

% nguồn: Ok_Toan11_TracNghiem_Lan2
\begin{ex}%[1D2N2-2]%[ma:7a72b014]
	Xác định số hạng đầu và công  sai của cấp số cộng $(u_n)$ với $u_n=2n+3$, $\forall n \in \mathbb{N}^*$.
	 \choice
		{$u_1=2$;  $d=3$}
		{$u_1=3$;  $d=2$}
		{\True $u_1=5$;  $d=2$}
		{$u_1=5$; $d=-3$} 
	\loigiai{ 
	Ta có ${u_1}=2\cdot 1+3=5$  và với mỗi $n \in \mathbb{N}^*$ bất kì, ta có $ d=u_{n+1}-u_n=\left[2\left(n+1\right)+3\right]-\left(2n+3\right)=2$.}
\end{ex}

% nguồn: Ok_Toan11_TracNghiem_Lan2
\begin{ex}%[1D2N2-2]%[ma:f1b53399]
	Trong các dãy số sau, dãy số nào \textbf{không phải} cấp  số cộng?
	\choice 
		{$1$;  $1$; $1$; $1$; $1$}
		{$-8$;  $-6$; $-4$; $-2$; $0$}
		{\True $3$; $1$; $-1$; $-2$; $-4$} 
		{$\dfrac{1}{2}$; $\dfrac{3}{2}$; $\dfrac{5}{2}$; $\dfrac{7}{2}$; $\dfrac{9}{2}$} 
	 \loigiai{
		\begin{enumerate}
			\item $1$; $1$; $1$; $1$; $1$ là cấp  số cộng với $u_1=1$, $d=0$.\\
			\item $-8$; $-6$; $-4$; $-2$; $0$ là cấp số cộng với $u_1=-8$, $d=2$.\\ 
			\item $3$; $1$;  $-1$; $-2$; $-4$ không là cấp số cộng vì $\left(-1\right)-1=-2\ne\left(-2\right)-\left(-1\right)=-1$.\\
			\item $\dfrac{1}{2}$; $\dfrac{3}{2}$; $\dfrac{5}{2}$;  $\dfrac{7}{2}$; $\dfrac{9}{2}$ là cấp số cộng có $u_1=\dfrac{1}{2}$, $d=1$.
		\end{enumerate}
	} 
\end{ex}

% nguồn: Ok_Toan11_TracNghiem_Lan2
\begin{ex}%[1D2N2-2]%[ma:76626fbb]
	Trong các dãy số sau, dãy nào là một cấp số cộng? 
	 \choice
		{$1$; $-3$; $-6$; $-9$; $-12$} 
		{\True $1$; $-3$;  $-7$; $-11$; $-15$}
		{$1$; $-3$; $-5$; $-7$; $-9$} 
		 {$1$; $-2$; $-4$; $-6$; $-8$}
	 \loigiai{
	Ta có dãy số $1$; $-3$; $-7$; $-11$; $-15$ là một cấp số  cộng có công sai $d=-4$.}
\end{ex}

% nguồn: Ok_Toan11_TracNghiem_Lan2
\begin{ex}%[1D2N2-2]%[ma:f81e6865]
	Cho cấp số cộng $1$, $8$, $15$, $22$, $29$, $\ldots $.  Công sai của cấp số cộng này là
	 \choice
		{$9$} 
		 {$1$}
		{$8$} 
		{\True  $7$}
	 \loigiai{
	 Công sai $d=u_2-u_1=8-1=7$.}
\end{ex}

% nguồn: Ok_Toan11_TracNghiem_Lan2
\begin{ex}%[1D2N2-2]%[ma:285eb91b]
	Trong các dãy số  sau đây, dãy số nào là cấp số cộng?
	\choice 
		 {$u_n=3^n$}
		 {$u_n=\left(-3\right)^{n+1}$}
		{\True  $u_n=3n+1$}
		{$u_n=2^{n+1}$} 
	\loigiai{ 
		Ta có 
		\begin{itemize}
			\item Xét $u_n=3^n$, $u_{n+1}-u_n=3^{n+1}-3^n=2\cdot3^n$ $\left(\forall n\in{\mathbb{N}^*}\right)$ nên $u_n=3^n$ không phải là  cấp số cộng.
			\item Xét $u_n=\left(-3\right)^{n+1}$, $u_{n+1}-u_n=\left(-3\right)^{n+1}-\left(-3\right)^n=-4\cdot \left(-3\right)^n$ $\left(\forall n\in{\mathbb{N}^*}\right)$ nên $u_n=\left(-3\right)^{n+1}$ không phải là cấp số cộng. 
			\item Xét  $u_n=3n+1$, $u_{n+1}-u_n =\left[3\left(n+1\right)+1\right]-\left(3n+1\right)=3$ $\left(\forall n\in{\mathbb{N}^*}\right)$ không đổi, nên $u_n=3n+1$ là cấp số cộng.
			\item Xét $u_n=2^{n+1}$, $u_{n+1}-u_n=2^{n+2}-2^{n+1}=2^{n+1}$ $\left(\forall n\in{\mathbb{N}^*}\right)$ nên $u_n=2^{n+1}$  không phải là cấp số cộng.
		\end{itemize}
	} 
\end{ex}

% nguồn: Ok_TracNghiemLop11_Lan1
\begin{ex}%[1D2N2-2]%[ma:50b5520e]
	Trong các dãy  số sau, dãy nào là cấp số cộng?
	 \choice
		{$2$;  $3$; $5$; $7$}
		 {$\dfrac{1}2$;$\dfrac{1}{4}$;$\dfrac{1}{8}$;$\dfrac{1}{16}$}
		 {$1$;$0$;$0$;$0$}
		 {$1$;$3$;$5$;$7$}
	\loigiai{Ta có cấp số cộng là dãy số mà mỗi số  hạng kể từ số hạng thứ hai bằng số đứng liên trước cộng với một số không đổi. Số không đổi gọi là công sai của cấp số cộng, kí hiệu là $d$.\\
		 Mà $3=1+2$; $5=3+ $; $7=5+2$.\\
		Vậy dãy số $1$; $3$; $5$; $7$ là cấp số cộng với công sai $d=2$. 
	 }
\end{ex}

% nguồn: Ok_TracNghiemLop11_Lan1
\begin{ex}%[1D2N2-2]%[ma:a000e2aa]
	Cho cấp số cộng $\left(u_n\right)$ có số hạng tổng quát là  $u_n=5n-1$. Tìm công sai $d$ của cấp số cộng.
	\choice 
		{\True $d=5$} 
		 {$d=6$}
		 {$d=-1$}
		 {$d=4$}
	\loigiai{ 
		Ta có $ d=u_n-u_{n-1}=5$. 
	 }
\end{ex}

% nguồn: Ok_TracNghiemLop11_Lan1
\begin{ex}%[1D2N2-2]%[ma:a19dfc02]
	Cho cấp số cộng có số hạng đầu $u_1=1$, công sai $d=-\dfrac{1}{3}$. Số hạng thứ $4$ của cấp số cộng  trên là
	 \choice
		{$-\dfrac{1}{3}$} 
		 {$\dfrac{2}{3}$}
		{$-2$} 
		{\True  $0$}
	\loigiai{ Áp dụng công thức tính số hạng tổng  quát của cấp số cộng $u_n=u_1+(n-1)\cdot d, \, \forall n\geq 1$.\\ Suy ra $u_4=u_1+3\cdot d \Rightarrow u_4=0$.
	 }
\end{ex}

% nguồn: Ok_TracNghiemLop11_Lan1
\begin{ex}%[1D2N2-2]%[ma:54b0cbe2]
	Cho cấp số cộng $(u_n)$ có $u_4=-12$ và $u_{14}=18$. Tìm số hạng đầu tiên $u_1$ và công sai $d$ của cấp số cộng đã  cho.
	\choice 
		{$u_1=-22$;  $d=3$}
		{\True $u_1=-21$;  $d=3$}
		{  $u_1=-21$; $d=-3$}
		{ $u_1=-20$;  $d=-3$}
	\loigiai{Ta  có
		$\heva{&-12=u_4=u_1+3d \\& 18=u_{14}=u_1+13d  \\
		}\Rightarrow \heva{ 
		 &u_1=-21 \\
	 & d=3.}$}
\end{ex}

% nguồn: Ok_TracNghiemLop11_Lan1
\begin{ex}%[1D2N2-2]%[ma:95d6db29]
	Viết ba số hạng  xen giữa các số $2$ và $22$ để được một cấp số cộng có năm số hạng.
	\choice 
		{\True $7$; $12$;  $17$}
		{$6$;  $10$; $14$}
		{$8$; $13$;  $18$}
		{$6$; $12$; $18$} 
	 \loigiai{
		Giữa $2$ và $22$ có thêm ba số hạng nữa lập thành cấp số cộng, xem như ta có một cấp số cộng có $5$ số hạng với $u_1=2 $; $u_5=22 $; ta cần tìm  $u_2$, $u_3$, $u_4 $.\\
		Ta có  $u_5=u_1+4d\Leftrightarrow d=\dfrac{u_5-u_1}{4}=\dfrac{22-2}{4}=5\Rightarrow \heva{&{u_2=u_1+d=7} \\
		 &{u_3=u_1+2d=12} \\
		&{u_4=u_1+3d=17} 
		 }$
	 }
\end{ex}

% nguồn: Ok_TracNghiemLop11_Lan1
\begin{ex}%[1D2N2-2]%[ma:20571e6b]
	Dãy nào sau đây là cấp  số cộng?
	\choice 
		{$1$; $5$;  $9$; $5$; $1$}
		{$1$; $3$;  $5$; $7$; $11$}
		{\True $2$; $5$;  $8$; $11$}
		{$3$; $7$; $3$; $7$} 
	 \loigiai{
		Dãy số $2$; $5$; $8$; $11$ là một cấp số cộng  với công sai $d=3$.
	} 
\end{ex}

% nguồn: Ok_TracNghiemLop11_Lan1
\begin{ex}%[1D2N2-2]%[ma:debf9e29]
	 Trong các dãy số sau, dãy số nào là cấp số cộng?
	 \choice
		{$1$;  $2$; $3$; $5$; $7$}
		 {\True $\dfrac{1}{3}$; $1$; $\dfrac{5}{3}$; $\dfrac{7}{3}$; $3$}
		{$\dfrac{1}{2}$;  $0$; $\dfrac{-1}{2}$; $1$; $\dfrac{-3}{2}$}
		{$0$;  $0$; $1$; $0$; $0$}
	\loigiai{ 
		 Dãy số $\dfrac{1}{3}$; $1$; $\dfrac{5}{3}$; $\dfrac{7}{3}$; $3$ là một cấp số cộng với công sai $d=\dfrac{2}{3}$.
	} 
\end{ex}

% nguồn: Ok_TracNghiemLop11_Lan1
\begin{ex}%[1D2N2-2]%[ma:8adf26b4]
	Dãy số nào sau đây không phải là cấp số  cộng?
	\choice 
		{$-\dfrac{2}{3};-\dfrac{1}{3};0;\dfrac{1}{3};\dfrac{2}{3};1;\dfrac{4}{3},\ldots$} 
		 {$15\sqrt{2};12\sqrt{2};9\sqrt{2};6\sqrt{2};\ldots$}
		{\True $\dfrac{4}{5};1;\dfrac{7}{5};\dfrac{9}{5};\dfrac{11}{5};\ldots$} 
		 {$\dfrac{1}{\sqrt{3}};\dfrac{2\sqrt{3}}{3};\sqrt{3};\dfrac{4\sqrt{3}}{3};\dfrac{5}{\sqrt{3}};\ldots$}
	\loigiai{ 
		Chỉ cần tồn  tại hai cặp số hạng liên tiếp của dãy số có hiệu khác nhau: $u_{m+1}-u_m \not=u_{k+1}-u_k$ thì ta kết luận ngay dãy số đó không phải là cấp số cộng.
		\begin{itemize}
			\item Xét  đáp án $-\dfrac{2}{3};-\dfrac{1}{3};0;\dfrac{1}{3};\dfrac{2}{3};1;\dfrac{4}{3};\ldots\mathop{\to}\limits_{} \dfrac{1}{3}=u_2-u_1=u_3-u_2=u_4-u_3=\cdots \mathop{\to}\limits_{}$ loại đáp án này.
			\item Xét đáp  án
			$15\sqrt{2};12\sqrt{2};9\sqrt{2};6\sqrt{2};\ldots\mathop{\to}\limits_{}-3\sqrt{3}=u_2-u_1=u_3-u_2=u_4-u_3=\cdots \mathop{\to}\limits_{}$  loại đáp án này.
			\item 
			Xét đáp án $\dfrac{4}{5};1;\dfrac{7}{5};\dfrac{9}{5};\dfrac{11}{5};\ldots\mathop{\to}\limits_{} \dfrac{1}{5}=u_2-u_1 \not=u_3-u_2=\dfrac{2}{5}  \mathop{\to}\limits_{}$ chọn đáp án này.
			\item Xét đáp án $\dfrac{1}{\sqrt{3}};\dfrac{2\sqrt{3}}{3};\sqrt{3};\dfrac{4\sqrt{3}}{3};\dfrac{5}{\sqrt{3}};\ldots\mathop{\to}\limits_{} \dfrac{\sqrt{3}}{3}=u_2-u_1=u_3-u_2=u_4-u_3 \mathop{\to}\limits_{}$ loại đáp án này. 
		\end{itemize}
	} 
\end{ex}

% nguồn: Ok_TracNghiemLop11_Lan1
\begin{ex}%[1D2N2-2]%[ma:ae8428bc]
	Tìm công thức số hạng tổng quát của cấp số cộng $(u_n)$ thỏa  mãn $\heva{& u_2-u_3+u_5=7 \\ & u_1+u_6=12.}$
	 \choice
		{$u_n=2n+3$} 
		{$u_n=2n+1$} 
		{$u_n=2n-3$} 
		 {\True $u_n=2n-1$}
	 \loigiai{
		Giả  sử dãy cấp số cộng $(u_n)$ có công sai là $ d$. \\
		Ta có 
		\allowdisplaybreaks 
		\begin{eqnarray*}
			\heva{& u_2-u_3+u_5=7 \\ & u_1+u_6=12 } &\Leftrightarrow& \heva{& ( u_1+d)-( u_1+2d)+( u_1+4d)=7  \\ & u_1+( u_1+5d)=12}\\
			&\Leftrightarrow& \heva{& u_1+3d=7  \\ & 2u_1+5d=12}\\
			&\Leftrightarrow& \heva{& u_1=1 \\  & d=2.}
		\end{eqnarray*}
		Số hạng tổng  quát của cấp số cộng $( u_n)$ là $$u_n=u_1+( n-1)d=1+( n-1)\cdot 2=2n-1.$$
		Vậy  $u_n=2n-1$.
	 }
\end{ex}

% nguồn: Ok_TracNghiemLop11_Lan1
\begin{ex}%[1D2N2-2]%[ma:06474243]
	Cho cấp số cộng  $(u_n)$ có $u_2=8$, $u_5=17$. Công sai $d$ bằng
	\choice 
		{$d=-3$} 
		{$d=-5$} 
		{\True $d=3$} 
		 {$d=5$}
	 \loigiai{
		 Ta có $ \heva{
		&u_2=8  \\
		&u_5=17 }  \Leftrightarrow \heva{
		&u_1+d=8\\ 
		 &u_1+4d=17 } \Leftrightarrow \heva{
		&u_1=5 \\ 
		&d=3.}$\\ 
	Vậy công sai  $d=3$.}
\end{ex}

% nguồn: Ok_TracNghiemLop11_Lan1
\begin{ex}%[1D2N2-2]%[ma:79e890a9]
	 Cho cấp số cộng $\left(u_n\right)$ với $u_1 = 2$ và $u_6 = 27$. Công sai $d$ của cấp số cộng đã cho là
	\choice 
		{$d  = -4$}
		{\True $d = 4$} 
		{$d  = 8$}
		{$d =  3$}
	 \loigiai{
		Ta  có công thức số hạng tổng quát của cấp số cộng là $u_n = u_1 + \left(n-1\right)d$.\\
		Khi đó $u_6 = 27 \Leftrightarrow  u_1 + 5d = 27 \Leftrightarrow d = 4$.
	 }
\end{ex}

% nguồn: Ok_TracNghiemLop11_Lan1
\begin{ex}%[1D2N2-2]%[ma:06bc7095]
	Cho cấp số cộng $(u_{n})$ có các  số hạng lần lượt là $\dfrac{7}{2};2;\dfrac{1}{2};-1;\ldots$. Công sai của dãy cấp số cộng đã cho là
	\choice 
		{$\dfrac{3}{2}$} 
		 {$\dfrac{1}{2}$}
		{$-\dfrac{1}{2}$} 
		{\True $-\dfrac{-3}{2}$} 
	\loigiai{ 
		 $u_2-u_1=2-\dfrac{7}{2}=-\dfrac{3}{2}=d$.
	} 
\end{ex}

% nguồn: Ok_Toan11_TracNghiem_Lan2
\begin{ex}%[1D2H2-2]%[ma:4196a022]
	Cho dãy số $(u_n)$ với $u_n=2n+1$. Số hạng thứ $2019$  của dãy là
	\choice 
		 {$4390$}
		 {\True $4039$}
		{$4930$} 
		{$4093$} 
	\loigiai{Số hạng thứ $2019$ của dãy là $u_{2019}=  2 \cdot 2019+1 =4039$.
	} 
\end{ex}

% nguồn: Ok_Toan11_TracNghiem_Lan2
\begin{ex}%[1D2H2-2]%[ma:d3e53bce]
	Tìm công  sai $d$ của cấp số cộng $(u_n)$, $n \in \mathbb{N}^*$ có $u_1=1$, $u_4=13$.
	 \choice
		 {$d=3$}
		{$d=\dfrac{1}{4}$} 
		{\True  $d=4$}
		{$d=\dfrac{1}{3}$} 
	\loigiai{Ta có $u_4=13 \Rightarrow u_1+3d=13 \Rightarrow 1+3d=13 \Rightarrow 3d=12 \Rightarrow  d=4$.
	 }
\end{ex}

% nguồn: Ok_Toan11_TracNghiem_Lan2
\begin{ex}%[1D2H2-2]%[ma:4968a0f3]
	Trong các dãy số nào  sau đây, dãy số nào là một cấp số cộng?
	\choice 
		 {$u_n=3^{n+1}$}
		{$u_n=\dfrac{2}{n+1}$} 
		{$u_n=\sqrt{n^2+1}$} 
		{\True  $u_n=\dfrac{5n-2}{3}$}
	\loigiai{Xét  dãy số $(u_n)$ với $u_n=\dfrac{5n-2}{3}$, $ \forall n \in \mathbb{N}^*$.\\
		Ta  có $u_{n+1}-u_n=\dfrac{5(n+1)-2}{3}-\dfrac{5n-2}{3}=\dfrac{5(n+1)-2-(5n-2)}{3}=\dfrac{5}{3}$, $ \forall n \in \mathbb{N}^*$.\\
		Vậy $(u_n)$ là một  cấp số cộng.
	 }
\end{ex}

% nguồn: Ok_Toan11_TracNghiem_Lan2
\begin{ex}%[1D2H2-2]%[ma:6f87f490]
	Cho $(u_n)$ là một cấp số cộng thỏa mãn $u_1+u_3=8$ và  $u_4=10$. Công sai của cấp số cộng đã cho bằng
	\choice 
		{\True $3$} 
		{$6$} 
		{$2$} 
		{$4$} 
	\loigiai{Ta có $\heva{&u_1+u_3=8\\&u_{4}=10} \Rightarrow \heva{&u_1+(u_1+2d)=8\\&u_1+3d=10}  \Rightarrow \heva{&2u_1+2d=8\\&u_1+3d=10} \Rightarrow \heva{&u_1=1\\&d=3.}$\\
		Vậy công sai  của câp số cộng đã cho là $d=3$.
	 }
\end{ex}

% nguồn: Ok_Toan11_TracNghiem_Lan2
\begin{ex}%[1D2H2-2]%[ma:c658fe90]
	Xác định số hạng đầu $u_1$ và công sai $d$ của cấp  số cộng $(u_n)$ có $u_9=5u_2$ và \\$u_{13}=2u_6+5$.
	 \choice
		{\True $u_1=3$ và $d=4$} 
		 {$u_1=3$ và $d=5$}
		{$u_1=4$  và $d=5$}
		{$u_1=4$ và $d=3$} 
	\loigiai{Ta có $\heva{&u_9=5u_2\\&u_{13}=2u_6+5} \Rightarrow  \heva{&u_1+8d=5(u_1+d)\\&u_1+12d=2(u_1+5d)+5} \Rightarrow \heva{&4u_1-3d=0\\&u_1-2d=-5} \Rightarrow \heva{&u_1=3\\&d=4.}$\\
		Vậy $u_1=3$  và $d=4$.
	} 
\end{ex}

% nguồn: Ok_Toan11_TracNghiem_Lan2
\begin{ex}%[1D2H2-2]%[ma:31ab8895]
	Cho cấp số  cộng $(u_n)$ có $u_1=-3$, $u_6=27$. Tính công sai $d$.
	\choice 
		{$d=7$} 
		{$d=5$} 
		{$d=8$} 
		 {\True $d=6$}
	\loigiai{Ta  có $u_6=u_1+5d \Rightarrow 27=(-3)+5d \Rightarrow 5d=30 \Rightarrow d=6$.
	 }
\end{ex}

% nguồn: Ok_Toan11_TracNghiem_Lan2
\begin{ex}%[1D2H2-2]%[ma:bed6d659]
	Trong các dãy số nào sau  đây, dãy số nào là cấp số cộng?
	 \choice
		{$u_n=\dfrac{1}{n}$} 
		{\True $u_n=u_{n-1}-2$, $\forall  n \ge 2$}
		 {$u_n=2^n-1$}
		{$u_n=2u_{n-1}$,  $\forall n \ge 2$}
	\loigiai{Xét dãy số  $(u_n)$ với $u_n=u_{n-1}-2$, $\forall n \ge 2$.\\
		Ta có $u_n-u_{n-1}=-2$, $\forall n  \ge 2$ nên $(u_n)$ là một cấp số cộng.
	 }
\end{ex}

% nguồn: Ok_Toan11_TracNghiem_Lan2
\begin{ex}%[1D2H2-2]%[ma:d2449a0e]
	Cho cấp số cộng  $(u_n)$ với $u_{17}=33$ và $u_{33}=65$ thì công sai bằng
	 \choice
		{$1$} 
		 {$3$}
		{$-2$} 
		{\True $2$} 
	\loigiai{Ta  có $\heva{&u_{17}=33\\&u_{33}=65} \Rightarrow \heva{&u_1+16d=33\\&u_1+32d=65} \Rightarrow \heva{&u_1=1\\&d=2.}$\\
	Vậy cấp số cộng trên có công  sai là $2$.}
\end{ex}

% nguồn: Ok_Toan11_TracNghiem_Lan2
\begin{ex}%[1D2H2-2]%[ma:f25a9321]
	Cho  dãy số $(u_n)$ với $u_n=1+2^n$. Khi đó số hạng $u_{2018}$ bằng
	 \choice
		{$2^{2018}$} 
		 {$2017+2^{2017}$}
		{\True$1+2^{2018}$} 
		 {$2018+2^{2018}$}
	\loigiai{Số hạng $u_{2018}$  là $u_{2018}=1+2^{2018}$.
	} 
\end{ex}

% nguồn: Ok_Toan11_TracNghiem_Lan2
\begin{ex}%[1D2H2-2]%[ma:8b2a66d0]
	Trong các dãy số nào sau  đây, dãy số nào là cấp số cộng?
	 \choice
		 {$u_n=3n^2+2017$}
		 {\True $u_n=3n+2018$}
		{$u_n=3^n$} 
		{$u_n=(-3)^{n+1}$} 
	\loigiai{Xét dãy số $(u_n)$  với $u_n=3n+2018$, $\forall n \in \mathbb{N}^*$.\\
		Ta có  $u_{n+1}-u_n=3(n+1)+2018-(3n+2018)=3n+3-3n=3$.\\
		 Vậy $(u_n)$ là một cấp số cộng.
	 }
\end{ex}

% nguồn: Ok_Toan11_TracNghiem_Lan2
\begin{ex}%[1D2H2-2]%[ma:dbc2f29d]
	Cho dãy số $(u_n)$ thỏa mãn $u_n=\dfrac{2n+1}{n}$. Tìm  số hạng thứ $10$ của dãy số đã cho.
	 \choice
		{\True $2{,}1$} 
		 {$2{,}2$}
		{$2{,}0$} 
		{$2{,}4$} 
	\loigiai{Số hạng thứ $10$ của dãy số  là $u_{10}=\dfrac{2 \cdot 10+1}{10}=\dfrac{21}{10}=2{,}1$.
	} 
\end{ex}

% nguồn: Ok_Toan11_TracNghiem_Lan2
\begin{ex}%[1D2H2-2]%[ma:025b01ae]
	Một cấp số cộng  gồm $5$ số hạng. Hiệu số hạng đầu và số hạng cuối bằng $20$. Tìm công sai $d$ của cấp số cộng đã cho
	\choice 
		{\True $d=-5$} 
		{$d=4$} 
		 {$d=-4$}
		{$d=5$} 
	\loigiai{Gọi  $u_1$, $u_2$, $u_3$, $u_4$, $u_5$ là cấp số cộng thỏa mãn yêu cầu đề bài.\\
		Ta có $u_1-u_5=20 \Rightarrow  u_1-(u_1+4d)=20 \Rightarrow -4d=20 \Rightarrow d=-5$.
	 }
\end{ex}

% nguồn: Ok_Toan11_TracNghiem_Lan2
\begin{ex}%[1D2H2-2]%[ma:f48911b2]
	Cho dãy số $(u_n)$ xác định  bởi $u_n=(-1)^n-\cos \left(n \pi\right)$. Giá trị $u_{99}$ bằng
	 \choice
		{\True $0$} 
		{$-1$} 
		{$1$} 
		 {$99$}
	\loigiai{Giá trị $u_{99}$ bằng  $u_{99}=(-1)^{99}-\cos \left(99 \pi\right)=-1-(-1)=0$.
	 }
\end{ex}

% nguồn: Ok_Toan11_TracNghiem_Lan2
\begin{ex}%[1D2H2-2]%[ma:dbd97bc7]
	Trong các dãy số nào sau đây, dãy số  nào là một cấp số cộng?
	 \choice
		{$u_n=n^2+1$, $\forall n \ge  1$}
		{$u_n=2^n$, $\forall n \ge 1$} 
		{$u_n=\sqrt{n+1}$, $\forall n  \ge 1$}
		{\True $u_n=2n-3$, $\forall n \ge  1$}
	\loigiai{Xét dãy số  $(u_n)$ với $u_n=2n-3$, $\forall n \ge 1$.\\
		Ta  có $u_{n+1}-u_n=2(n+1)-3-(2n-3)=2n+2-2n=2$.\\
		Vậy $(u_n)$ là một cấp số  cộng.
	 }
\end{ex}

% nguồn: Ok_Toan11_TracNghiem_Lan2
\begin{ex}%[1D2H2-2]%[ma:3df83cfa]
	Cho cấp số cộng $(u_n)$ có số hạng tổng quát là $u_n=3n-2$. Tìm công sai $d$  của cấp số cộng.
	 \choice
		{\True$d=3$} 
		 {$d=2$}
		{$d=-2$} 
		{$d=-3$} 
	 \loigiai{Ta có $d=u_{n+1}-u_n=3(n+1)-2-(3n-2)=3n+3-2-3n+2=3$.
	} 
\end{ex}

% nguồn: Ok_TracNghiemLop11_Lan1
\begin{ex}%[1D2H2-2]%[ma:349be0a2]
	Cho cấp số  cộng $\left(u_n\right)$ có $u_1=2024$ và công sai $ d=-20$. Tìm số hạng $u_{100}$.
	\choice 
		{$u_{100}=2024\cdot  \left(-20\right)^{99}$}
		{\True $u_{100}=44$} 
		 {$u_{100}=924$}
		{$u_{100}=946$} 
	\loigiai{ 
	 Ta có $u_{100}=u_1+99\cdot d=2024+99\cdot\left(-20\right)=44$.}
\end{ex}

% nguồn: Ok_TracNghiemLop11_Lan1
\begin{ex}%[1D2H2-2]%[ma:3b630d2d]
	Tìm công  sai $d$ của cấp số cộng hữu hạn biết số hạng đầu $u_1=10$ và số hạng cuối $u_{21}=50$.
	\choice 
		{$d=3$} 
		{\True  $d=2$}
		 {$d=4$}
		 {$d=-2$}
	\loigiai{ 
		Ta có $u_{21}=u_1+20d$ $\Rightarrow d=\dfrac{u_{21}-u_1}{20}$  $=\dfrac{50-10}{20}=2$.
	 }
\end{ex}

% nguồn: Ok_TracNghiemLop11_Lan1
\begin{ex}%[1D2H2-2]%[ma:c2227e61]
	Dãy số  nào sau đây là một cấp số cộng?
	\choice 
		{\True $(u_n)\colon \heva{ 
		 &u_1=2 \\
		&{u_{n+1}}=u_n+3,\,\,\forall  n\ge 1}$}
		{$(u_n)\colon  \heva{
		 &u_1=3 \\
		&{u_{n+1}}=2u_n-4,\,\,\forall n\ge 1}$} 
		{$(u_n)\colon  1$; $-3$; $6$; $7$; $15$; $\ldots$}
		{$(u_n)\colon -1$;  $1$; $-1$; $1$; $-1$; $\ldots$}
	 \loigiai{
		 Dãy số $(u_n)\colon \heva{
		&u_1=2 \\ 
	&{u_{n+1}}=u_n+3,\,\,\forall n\ge 1 } $ thỏa  $u_{n+1}-u_n=3$ với mọi $n\ge 1$ nên là cấp số cộng.}
\end{ex}

% nguồn: Ok_TracNghiemLop11_Lan1
\begin{ex}%[1D2H2-2]%[ma:82576a8a]
	Cho cấp số cộng $\left(u_n \right)$ thỏa mãn $u_1=4$, $u_3=10$. Công sai của cấp số cộng  bằng
	 \choice
		{$6$} 
		 {$-6$}
		{\True  $3$}
		{$-3$} 
	 \loigiai{
		Gọi $d$  là công sai của cấp số cộng $\left(u_n \right)$.\\
		Ta có $u_1=4,\;  u_3=10$ suy ra $\heva{&{u_1=4} \\
		 &{u_1+2d=10}
		&} \Leftrightarrow \heva{&{u_1=4} \\ 
		 &{d=3.}}$\\
		Vậy  công sai của cấp số cộng $\left(u_n \right)$ là $d=3$.
	} 
\end{ex}

% nguồn: Ok_Toan11_TracNghiem_Lan2
\begin{ex}%[1D2V2-2]%[ma:b69614b3]
	Cho  $(u_n)$ là một cấp số cộng thỏa mãn $u_1+{u_3}=8$ và ${u_4}=10$. Công sai của cấp số cộng đã cho bằng
	 \choice
		{\True  $3$}
		 {$6$}
		 {$2$}
		 {$4$}
	 \loigiai{
		Ta có $ u_3=u_1+2d, \,u_4 =u_1+3d. $\\ 
		Theo giả  thiết
		\[ 
		\heva{ 
		& u_1+u_3=8\\ 
		 & u_4=10}
		\Leftrightarrow 
		 \heva{
		 & u_1+(u_1+2d)=8\\
		& u_1+3d=10} 
		 \Leftrightarrow
		\heva{ 
		 & 2u_1+2d=8\\
		& u_1+3d=10.} 
		\] 
		Giải  hệ
		\[ 
		\heva{ 
		 & u_1=1\\
		 & d=3.}
		 \]
		Vậy  công sai của cấp số cộng là $d=3$.
	 }
\end{ex}

% nguồn: Ok_TracNghiemLop11_Lan1
\begin{ex}%[1D2V2-2]%[ma:b7ff6707]
	Cho cấp số cộng  $\left(u_n \right)$ có $\heva{&{u_2+u_5=42} \\
	 &{u_4+u_9=66}
	}$. Tổng $346$ số hạng đầu tiên  của cấp số cộng trên là
	\choice 
		 {$S_{346}=422554$}
		 {\True $S_{346}=242546$}
		 {$S_{346}=156224$}
		 {$S_{346}=-203558$}
	\loigiai{ 
		Ta có $\heva{&{u_2+u_5=42}  \\
		&{u_4+u_9=66} 
		} \Leftrightarrow \heva{&{2u_1+5d=42}  \\
		 &{2u_1+11d=66}
		 } \Leftrightarrow \heva{&{u_1=11} \\
		 &{d=4}.}$\\
		Và  $S_n=nu_1+\dfrac{n\left(n-1\right)d}{2} \Rightarrow S_{346}=346\cdot11+\dfrac{346\cdot345\cdot4}{2}=242546$.
	 }
\end{ex}


%%%%%%%%%%%%%%%% ĐÚNG/SAI %%%%%%%%%%%%%%%%

% nguồn: Ok_dungsai_lop11_Lan1
\begin{ex}%[1D2N2-2]%[ma:7464f84e]
 Cho dãy số hữu  hạn gồm các số hạng $-1$; $2$; $5$; $8$; $11$; $14$; $17$ . Khi đó
  \choiceTF
  {Dãy số đã cho là không phải  cấp số cộng}
  {\True Số hạng  $u_1=-1$}
  {Nếu dãy số đã cho là một cấp số cộng thì công sai của cấp số  cộng là $d=2$}
  {\True Tổng tất cả  số hạng của dãy số bằng $56$}
 \loigiai{ 
  \begin{itemchoice}
    \itemch Đặt $u_1=-1;u_2=2;u_3=5;u_4=8;u_5=11;u_6=14;u_7=17$.\\
    Ta có $u_2-u_1=u_3-u_2=u_4-u_3=u_5-u_4=u_6-u_5=u_7-u_6=3$.\\
   Vậy dãy số hưu hạn  đã cho là một cấp số cộng.
   \itemch  $u_1=-1$.
   \itemch Công sai cấp số cộng là $d=3$. 
    \itemch Với $u_1=-1,n=7,d=3$ thì $S_n=\dfrac{n\left[2u_1+(n-1)d\right]}{2}=\dfrac{7[2(-1)+6\cdot3]}{2}=56$.
  \end{itemchoice}
 } 
\end{ex}

% nguồn: Ok_dungsai_lop11_Lan2
\begin{ex}%[1D2N2-2]%[ma:534292bf]
	Cho các dãy số sau:  $a_n=4(n+1)-3$, $b_n=4n^2$, $z_n=\dfrac{2-3n}{4}(n\ge 1)$.
	\choiceTF 
		{\True Hai số hạng đầu của dãy $a_n$ là $a_1=5$,  $a_2=9$}
		{Dãy $b_n$ là một cấp số cộng} 
		{\True Dãy $z_n$  là một cấp số cộng với công sai $d=-\dfrac{3}{4}$}
		{\True Có $2$ dãy số là cấp số  cộng}
	 \loigiai{
		\begin{itemchoice}
			\itemch {\bf  Đúng}.\\ Thay $n=1\Rightarrow a_1=5$ và $n=2\Rightarrow {a_n}=9$.
			\itemch {\bf Sai}.\\ Vì $b_1=4$,  $b_2=16$, $b_3=36$. Dễ thấy $16-4=12$, $36-16=20$. Do đó $b_n$ không là $1$ cấp số cộng.
			\itemch {\bf Đúng}.\\ Vì $z_{n+1}-z_n=\dfrac{2-3(n+1)}{4}-\dfrac{2-3n}{4}=\dfrac{-3n-1}{4}-\dfrac{2-3n}{4}=-\dfrac{3}{4}.$\\ 
			Vậy đây là cấp số cộng với công sai $d=-\dfrac{3}{4}.$ 
			\itemch {\bf Đúng}.\\ Vì $a_{n+1}-a_n=4(n+2)-3-\left[4(n+1)-3\right]=4$. Vậy $a_n$  là $1$ cấp số cộng. Do b) và c) nên kết luận có $2$ cấp số cộng là $a_n$ và $z_n$.
		\end{itemchoice}
	 }
\end{ex}

% nguồn: Ok_dungsai_lop11_Lan2
\begin{ex}%[1D2N2-2]%[ma:a8b830d1]
	Cho cấp số cộng  $(u_n)$ có $u_6=20$, $u_8=30$.
	 \choiceTF
		{\True Số hạng $u_7=25$} 
		{\True Số hạng  $u_3=10$}
		{Công sai của cấp số cộng  bằng $d=-5$}
		{\True Số hạng $u_{2n}\ge 0, \forall n\ge 1$} 
	 \loigiai{Ta có
		\begin{eqnarray*}
			&  & \heva{&u_6=0\\&u_8=30}\\
			&\Leftrightarrow & \heva{&u_1+5d=0\\&u_1+7d=30}\\ 
			&\Leftrightarrow  & \heva{&u_1=-5\\&d=5.}
		\end{eqnarray*}
		\begin{itemchoice}
			\itemch {\bf Đúng}.\\  Vì số hạng $u_7=\dfrac{u_6+u_8}{2}=25$.
			\itemch {\bf Đúng}.\\ Vì  số hạng $u_3=u_1+2d=10$.
			\itemch  {\bf Sai}.\\ Vì $d=5$.
			\itemch {\bf Đúng}.\\ Vì $u_n=u_1+(n-1)d=-5+(n-1)\cdot 5=5n-10\Rightarrow {u_n}=10n-10\ge 0,  \forall n\ge 1$.
		\end{itemchoice}
	 }
\end{ex}

% nguồn: Ok_dungsai_lop11_Lan2
\begin{ex}%[1D2N2-2]%[ma:0971bc26]
	Cho dãy số hữu hạn gồm các số hạng:  $-1$; $2$; $5$; $8$; $11$; $14$; $17$. Khi đó
	\choiceTF 
		{Dãy số đã cho  là không phải cấp số cộng}
		{\True Số hạng  $u_1=-1$, $n=7$, $d=3$}
		{Nếu dãy số đã cho là  một cấp số cộng thì công sai của cấp số cộng là $d=2$}
		{\True Tổng tất  cả số hạng của dãy số bằng $56$}
	 \loigiai{
		Đặt: $u_1=-1$; $u_2=2$; $u_3=5$; $u_4=8$;  $u_5=11$; $u_6=14$; $u_7=17$.\\
		Ta có  $u_2-u_1=u_3-u_2=u_4-u_3=u_5-u_4=u_6-u_5=u_7-u_6=3$.\\
		Vậy dãy số hữu hạn đã cho là một cấp số cộng.\\ 
		Công sai cấp số  cộng là $d=3$.\\
		Với $\heva{&u_1=-1 \\&n=7 \\&d=3}$ thì $S_n=\dfrac{n\big[2u_1+(n-1) d\big]}{2}=\dfrac{7\big[2(-1)+6\cdot  3\big]}{2}=56$.
		\begin{itemchoice}
			\itemch {\bf  Sai}.\\
			\itemch {\bf Đúng}.\\ 
			 \itemch {\bf Sai}.\\
			\itemch {\bf  Đúng}.\\
		\end{itemchoice}
	 }
\end{ex}

% nguồn: Ok_dungsai_lop11_Lan2
\begin{ex}%[1D2N2-2]%[ma:d1ae8a32]
	Cho dãy số: $-2$; $1$;  $4$; $7$; $10$.
	 \choiceTF
		{Đây không  phải là cấp số cộng}
		{Đây là một cấp số  cộng hữu hạn}
		{Đây  là một cấp số cộng với công sai $d=-2$}
		{\True Số hạng tiếp  theo của dãy số là $13$}
	 \loigiai{
		Dễ nhận thấy $1=-2+3$, $4=1+3$, $7=4+3$, $10=7+3$. Do đó dãy $-2$; $1$; $4$; $7$;  $10$ là một cấp số cộng với số hạng đầu $u_1=-2$ và công sai $d=3$. Vậy
		\begin{itemchoice}
			\itemch {\bf  Sai}.\\
			\itemch  {\bf Sai}.\\
			\itemch {\bf Sai}.\\ Vì công sai $d=3$. 
			\itemch {\bf  Đúng}.\\ Vì số hạng tiếp theo $10+3=13$.
		\end{itemchoice}
	 }
\end{ex}

% nguồn: Ok_dungsai_lop11_Lan1
\begin{ex}%[1D2H2-2]%[ma:5bc5e28f]
 Cho cấp số cộng $1$; $3$; $5$; $7$; $9$; $11$. Các mệnh  đề sau đúng hay sai?
 \choiceTF 
  {\True Số hạng đầu  $u_1$ bằng $1$}
  {Công sai của cấp số cộng bằng $3$} 
  {Số hạng thứ tư  bằng $4$}
  {\True Tổng ba số hạng đầu bằng  $9$}
 \loigiai{ 
  \begin{itemchoice}
   \itemch Ta có số hạng đầu $u_1$  bằng $1$.
   \itemch Công sai  của cấp số cộng bằng $3-1=2$.
   \itemch Số hạng  thứ tư bằng $7$.
   \itemch Tổng ba số hạng đầu bằng  $1+3+5=9$.
  \end{itemchoice}
 } 
\end{ex}

% nguồn: Ok_dungsai_lop11_Lan1
\begin{ex}%[1D2H2-2]%[ma:e97de46f]
  Cho cấp số cộng $\left(u_n \right)$ có $u_1=5$ và $d=-7$. Khi đó
 \choiceTF 
  {\True  $u_{11}=-65$}
  {$u_5+u_7=-50$} 
  {\True Số $-849$ là số hạng thứ  $123$ của cấp số cộng}
  {\True Số $-114$ là số hạng  thứ $18$ của cấp số cộng}
  \loigiai{
  Công thức số hạng tổng quát  của cấp số cộng là $u_n=u_1+(n-1)d=5+(n-1)\cdot (-7)=-7n+12$.
  \begin{itemchoice}
   \itemch  Ta có $u_{11}=-7\cdot11+12=-65$.
   \itemch  $u_5+u_7=-60$.
    \itemch Ta có $-849=-7n+12\Rightarrow n=123$.
   \itemch Ta có $-114=-7n+12\Rightarrow  n=18$.
  \end{itemchoice}
 } 
\end{ex}

% nguồn: Ok_dungsai_lop11_Lan1
\begin{ex}%[1D2H2-2]%[ma:d2a9b94a]
  Cho cấp số cộng $\left(u_n \right)$ có công sai $d < 0$ thoả mãn $\heva{&{u_1+u_7=26} \\
 &{u_2^2+u_6^2=466} 
  }$. Khi đó
 \choiceTF 
  {\True Số  hạng $u_1=25$}
  {Công  sai $d=-3$}
   {\True Số hạng $u_{10}=-11$}
   {\True Số hạng $u_{2024}=-8067$}
  \loigiai{Ta có
  $\heva{&u_1+u_7=26 \\&u_2^2+u_6^2=466}\Leftrightarrow \heva{&2u_1+6d=26  \\
   &\left(u_1+d\right)^2+\left(u_1+5d\right)^2=466}\Leftrightarrow \heva{u_1=13-3d\qquad \qquad \qquad (1) \\
  &\left(u_1+d\right)^2+\left(u_1+5d\right)^2=466\qquad (2).}$\\ 
  Thay (1) vào  (2), ta được: $(13-2d)^2+(13+2d)^2=466\Leftrightarrow 8d^2+338=466$
  \[\Leftrightarrow  \hoac{&{d=4} \\
  &{d=-4}} 
   \]
  Vì $d < 0$ nên ta nhận $d=-4$,  khi đó $u_1=25$.
  \begin{itemchoice}
   \itemch  $u_1=25$.
   \itemch  $d=-4$.
    \itemch $u_{10}=u_1+9d=-11$.
   \itemch Ta có $u_n=u_1+(n-1)d=25+(n-1)(-4)=29-4n$. Suy ra  $u_{2024}=-8067$.
  \end{itemchoice}
 } 
\end{ex}

% nguồn: Ok_dungsai_lop11_Lan1
\begin{ex}%[1D2H2-2]%[ma:319ee537]
 Cho cấp  số cộng $3$; $6$; $9$; $12$;\ldots. Các mệnh đề sau đúng hay sai?
  \choiceTF
  {\True Số hạng đầu $u_1$  bằng $3$}
  {Công sai của cấp số cộng  bằng $6$}
  {Số  hạng thứ năm bằng $24$}
  {\True Số hạng  tổng quát là $u_n=3n$}
 \loigiai{ 
  \begin{itemchoice}
   \itemch Ta có số hạng đầu  $u_1$ bằng $3$.
   \itemch Công sai của cấp  số cộng bằng $6-3=3$.
   \itemch Ta  có $u_4=12$ nên $u_5=u_4+d=12+3=15$.
   \itemch Số hạng tổng quát  $u_n=u_1+(n-1)d=3+(n-1)\cdot3=3n$.
  \end{itemchoice}
 } 
\end{ex}

% nguồn: Ok_dungsai_lop11_Lan2
\begin{ex}%[1D2H2-2]%[ma:c28c6ebf]
	Cho  dãy số $(u_n)$ xác định bởi: $u_1=2$ và $u_{n+1}=5+u_n$, $\forall n\in{\mathbb{N}^{*}}$.
	 \choiceTF
		{Số  hạng $u_2=6$}
		 {\True Tỉ số $u_{2\,024}-u_{2\,023}=5$}
		{\True $(u_n)$ là cấp số cộng với công sai của cấp số cộng bằng $d=5$} 
		{Tổng $S=\dfrac{1}{u_1u_2}+\dfrac{1}{u_2u_3}+\ldots +\dfrac{1}{u_{48}u_{49}}+\dfrac{1}{u_{49}u_{50}}$ lớn hơn  $1$}
	 \loigiai{
		\begin{itemchoice}
			\itemch {\bf Sai}.\\ Vì số  hạng $u_2=u_1+5=7$.
			\itemch {\bf Đúng}.\\ Vì $u_{n+1}=5+u_n$, $\forall n\in{\mathbb{N}^{*}}\Rightarrow {u_{n+1}}-u_n=5$ nên $u_{2\,024}-u_{2\,023}=5$. 
			\itemch {\bf Đúng}.\\ Có $u_{n+1}=5+u_n$, $\forall n\in{\mathbb{N}^{*}}\Rightarrow u_{n+1}-u_n=5$ nên  $(u_n)$ là cấp số cộng với công sai của cấp số cộng bằng $d=5$.
			\itemch  {\bf Sai}.\\ Ta có
			\begin{align*}
				d\cdot S=\dfrac{d}{u_1u_2}+\dfrac{d}{u_2u_3}+\ldots +\dfrac{d}{u_{48}u_{49}}+\dfrac{d}{u_{49}u_{50}}=\dfrac{1}{u_1}-\dfrac{1}{u_2}+\dfrac{1}{u_2}-\dfrac{1}{u_3}+\ldots  +\dfrac{1}{u_{49}}-\dfrac{1}{u_{50}}=\dfrac{1}{u_1}-\dfrac{1}{u_{50}}.
			\end{align*}
			Suy ra  $S=\dfrac{1}{d}\left(\dfrac{1}{u_1}-\dfrac{1}{u_{50}}\right)$.\\
			Với $u_1=1$; $d=5\Rightarrow {u_{50}}=u_1+49d=246$ ta được  $S=\dfrac{1}{5}\left(1-\dfrac{1}{246}\right)=\dfrac{49}{246} < 1$.
	\end{itemchoice}}
\end{ex}

% nguồn: Ok_dungsai_lop11_Lan2
\begin{ex}%[1D2H2-2]%[ma:a445abaf]
	Cho các dãy  số có số hạng tổng quát $a_n=4n-3$, $b_n=\dfrac{2-3n}{4}$, $c_n=n^2$. Khi đó
	 \choiceTF
		{\True $\left(a_n\right)$ là một cấp số  cộng với số hạng đầu $a_1=1$}
		{\True  $\left(a_n\right)$ là một cấp số cộng với công sai $d=4$}
		{$\left(b_n\right)$  là một cấp số cộng với số hạng đầu $b_1=-\dfrac{1}{4}$ và công sai $d=\dfrac{3}{4}$}
		{$\left(c_n\right)$ là  một cấp số cộng với công sai $d=2$}
	\loigiai{ 
		\begin{itemchoice}
			\itemch  \textbf{Đúng.}\\
			Ta có  $a_{n+1}-a_n=4(n+1)-3-(4n-3)=4,\forall n\ge 1$.\\
			Do đó $\left(a_n\right)$ là một cấp số cộng với số  hạng đầu $a_1=4\cdot 1-3=1$ và công sai $d=4$.
			\itemch  \textbf{Đúng.}
			\itemch  \textbf{Sai.}\\
			Ta  có $b_{n+1}-b_n=\dfrac{2-3(n+1)}{4}-\dfrac{2-3n}{4}=\dfrac{2-3n-3-2+3n}{4}=-\dfrac{3}{4},\forall n\ge 1$.\\
			Suy ra $\left(b_n\right)$ là một cấp  số cộng với số hạng đầu $b_1=\dfrac{2-3.1}{4}=-\dfrac{1}{4}$ và công sai $d=-\dfrac{3}{4}$.
			 \itemch \textbf{Sai.}\\
			Ta có $c_{n+1}-c_n=(n+1)^2-n^2=2n+1$ (phụ thuộc  vào giá trị của $n$).\\
			Suy ra $\left(c_n\right)$  không phải là một cấp số cộng.
		\end{itemchoice}
	} 
\end{ex}

% nguồn: Ok_dungsai_lop11_Lan2
\begin{ex}%[1D2H2-2]%[ma:06da60a3]
	Cho cấp số cộng $(u_n)$ biết  $u_n=n+5$.
	 \choiceTF
		{\True Số hạng đầu của cấp số cộng là  $u_1=6$}
		{Công sai của cấp số cộng bằng  $5$}
		{\True Tổng $T=u_2+u_4+u_6+\ldots  +u_{2\,024}=4\,060\,216$}
		{Tổng $S=\dfrac{1}{\sqrt{u_1}+\sqrt{u_2}}+\dfrac{1}{\sqrt{u_2}+\sqrt{u_3}}+  \ldots +\dfrac{1}{\sqrt{u_{2\,019}}+\sqrt{u_{2\,020}}}=\dfrac{2\,019}{45}$}
	 \loigiai{
		\begin{itemchoice}
			\itemch {\bf Đúng}.\\ Vì  thay $n=1$ ta được $u_1=6$.
			\itemch {\bf Sai}.\\ Vì $u_1=6$; $u_2=7\Rightarrow d=u_2-u_1=1$. 
			\itemch {\bf Đúng}.\\ Ta có $u_2=7$; $u_4=9$; $u_6=11$; $\ldots$; $u_{2\,024}=2\,029$. Dãy số $7$; $9$; $11$; $\ldots$; $2\,024$ lập thành một cấp số cộng với $u'_1=7$; $d'=2$ và $2\,029$ là  số hạng thứ $2\,012$ của cấp số cộng này.\\
			Vậy  $T=u_2+u_4+u_6+\ldots +u_{2\,024}=7\cdot 2\,012+\dfrac{2\,012\cdot 2\,011}{2}\cdot 2=4\,060\,216$.
			\itemch {\bf Sai}.\\  Ta có $u_{2\,020}=2\,025$. Suy ra
			\begin{align*}S&=\dfrac{1}{\sqrt{u_1}+\sqrt{u_2}}+\dfrac{1}{\sqrt{u_2}+\sqrt{u_3}}+\ldots +\dfrac{1}{\sqrt{u_{2\,019}}+\sqrt{u_{2\,020}}}\\
				&=\dfrac{\sqrt{u_2}-\sqrt{u_1}}{u_2-u_1}+\dfrac{\sqrt{u_3}-\sqrt{u_2}}{u_3-u_2}+\ldots +\dfrac{\sqrt{u_{2\,020}}-\sqrt{u_{2\,019}}}{u_{2\,020}-u_{2\,019}}\\ 
				&=\sqrt{u_2}-\sqrt{u_1}+\sqrt{u_3}-\sqrt{u_2}+\ldots  +\sqrt{u_{2\,020}}-\sqrt{u_{2\,019}}\\
				&=\sqrt{u_{2\,020}}-\sqrt{u_1}=\dfrac{u_{2\,020}-u_1}{\sqrt{u_{2\,020}}+\sqrt{u_1}}=\dfrac{2\,019}{45+\sqrt{6}}. 
			\end{align*}
		\end{itemchoice}
	 }
\end{ex}

% nguồn: Ok_dungsai_lop11_Lan2
\begin{ex}%[1D2H2-2]%[ma:33e49654]
	Cho một cấp số cộng $(u_n)$  có $u_1=1$ và $d=5$.
	\choiceTF 
		{Số hạng thứ hai của cấp số cộng là  $u_2=5$}
		{\True Tổng $6$ số hạng đầu  của cấp số cộng là $S_6=\dfrac{6\left(u_1+u_6\right)}{2}$}
		{Tổng  của $10$ số hạng đầu của cấp số cộng là $460$}
		{\True Tổng  $S=\dfrac{1}{u_1^{}u_2}+\dfrac{1}{u_2u_3}+\ldots +\dfrac{1}{u_{49}u_{50}}=\dfrac{49}{246}$}
	 \loigiai{
		\begin{itemchoice}
			\itemch {\bf  Sai}.\\ $u_2=1+5=6$.
			 \itemch {\bf Đúng}.\\ Theo công thức $S_n=\dfrac{\left(u_1+u_n\right)n}{2}.$
			\itemch {\bf Sai}.\\ $S_{10}=10u_1+\dfrac{10\cdot 9}{2}\cdot d=10\cdot 1+\dfrac{10\cdot  9}{2}\cdot 5=235$.
			\itemch {\bf  Đúng}.\\ Ta có
			\begin{align*}5S&=\dfrac{5}{u_1u_2}+\dfrac{5}{u_2u_3}+\ldots +\dfrac{5}{u_{49}u_{50}}\\
				 &=\dfrac{u_2-u_1}{u_1u_2}+\dfrac{u_3-u_2}{u_2u_3}+\ldots +\dfrac{u_{50}-u_{49}}{u_{49}u_{50}}\\
				&=\dfrac{1}{u_1}-\dfrac{1}{u_2}+\dfrac{1}{u_2}-\dfrac{1}{u_3}+\ldots +\dfrac{1}{u_{48}}-\dfrac{1}{u_{49}}+\dfrac{1}{u_{49}}-\dfrac{1}{u_{50}}\\ 
				&=\dfrac{1}{u_1}-\dfrac{1}{u_{50}}=\dfrac{1}{u_1}-\dfrac{1}{u_1+49d}=\dfrac{245}{246}. 
			\end{align*}
			Suy  ra $S=\dfrac{49}{246}$.
	\end{itemchoice}}
\end{ex}

% nguồn: Ok_dungsai_lop11_Lan1
\begin{ex}%[1D2V2-2]%[ma:a61e08f0]
 Xác định tính đúng, sai  của các khẳng định sau:
 \choiceTF 
  {\True Dãy số $\left(u_n \right)$ với $\dfrac{-2}{3}$; $\dfrac{-1}{3}$;  $0$; $\dfrac{1}{3}$; $\dfrac{2}{3}$; $1$; $\dfrac{4}{3}$ là cấp số cộng với $u_1=\dfrac{-2}{3}$; $d=\dfrac{1}{3}$}
   {\True Dãy số $\left(u_n \right)$ với $u_n=7-3n$ là cấp số cộng với $u_1=4$; $d=-3$}
  {Dãy số $\left(u_n \right)$ với $u_n=n^2+n+1$ là  cấp số cộng với $u_1=3$; $d=1$}
  {\True Dãy số $\left(u_n \right)$ với $u_n=(-1)^n+3n$ không  là cấp số cộng}
 \loigiai{ 
  \begin{itemchoice}
   \itemch Dãy số $\left(u_n \right)$  với $\dfrac{-2}{3};\dfrac{-1}{3};0;\dfrac{1}{3};\dfrac{2}{3};1;\dfrac{4}{3}$.\\
    Ta thấy $u_2-u_1=u_3-u_2=u_4-u_3=\ldots=\dfrac{1}{3}$.\\
   Vậy $\left(u_n \right)$ là cấp số cộng với $u_1=\dfrac{-2}{3};d=\dfrac{1}{3}$. 
   \itemch  Dãy số $\left(u_n \right)$ với $u_n=7-3n$.\\
   Ta có  $u_{n+1}-u_n=[7-3(n+1)]-(7-3n)=-3$.\\
   Vậy $\left(u_n \right)$ là cấp số  cộng với $u_1=7-3\cdot1=4$; $d=-3$.
   \itemch Dãy số $\left(u_n  \right)$ với $u_n=n^2+n+1$.\\
   Ta có $u_{n+1}-u_n=(n+1)^2+(n+1)+1-\left(n^2+n+1\right)=2n+2$ phụ thuộc vào $n$.\\ 
   Vậy $\left(u_n  \right)$ không là cấp số cộng.
   \itemch Dãy  số $\left(u_n \right)$ với $u_n=(-1)^n+3n$.\\
   Ta có  $u_{n+1}-u_n=(-1)^{n+1}+3(n+1)-\left[(-1)^n+3n\right]=-(-1)^n+3-(-1)^n=3-2(-1)^n$ phụ thuộc vào $n$.\\
   Vậy $\left(u_n \right)$ không  là cấp số cộng.
  \end{itemchoice}
 } 
\end{ex}

% nguồn: Ok_dungsai_lop11_Lan2
\begin{ex}%[1D2V2-2]%[1D2H2-6]%[ma:a0563a41]
	Một  cấp số cộng có số hạng đầu $u_1=1$, công sai $d=3$.
	 \choiceTF
		{\True  Số hạng $u_2=4$}
		{\True Số hạng tổng quát $u_n=3n-2, (n\ge 1)$} 
		{Biết $S_n=14\,950$ thì $n=120$} 
		{$S=\dfrac{1}{u_2\sqrt{u_1}+u_1\sqrt{u_2}}+\dfrac{1}{u_3\sqrt{u_2}+u_2\sqrt{u_3}}+\ldots  +\dfrac{1}{u_{2\,024}\sqrt{u_{2\,023}}+u_{2\,023}\sqrt{u_{2\,024}}}=\dfrac{2\,023}{2\,024}$}
	 \loigiai{
		\begin{itemchoice}
			\itemch {\bf Đúng}.\\ Vì  số hạng $u_2=u_1+d=4$.
			\itemch {\bf Đúng}.\\ Vì số hạng tổng quát $u_n=u_1+(n-1)d=1+(n-1)3=3n-2$. 
			\itemch {\bf Sai}.\\ Áp dụng  công thức
			\begin{eqnarray*}
				& & S_n=nu_1+\dfrac{n\cdot  (n-1)}{2}d\\
				&\Leftrightarrow & n+\dfrac{n(n-1)3}{2}=14\,950\\ 
				&\Leftrightarrow &  n=100.
			\end{eqnarray*}
			\itemch {\bf  Sai}.\\ Ta có $u_{2\,024}=3\cdot 2\,024-2=6\,070$.
			Khi  đó
			\begin{align*}
				\dfrac{1}{u_{k+1}\sqrt{u_k}+u_k\sqrt{u_{k+1}}}=\dfrac{1}{\sqrt{u_k}\cdot \sqrt{u_{k+1}}\cdot  \left(\sqrt{u_k}+\sqrt{u_{k+1}}\right)}=\dfrac{1}{d}\cdot \dfrac{\sqrt{u_{k+1}}-\sqrt{u_k}}{\sqrt{u_k}\cdot \sqrt{u_{k+1}}}=\dfrac{1}{d}\cdot \left(\dfrac{1}{\sqrt{u_k}}-\dfrac{1}{\sqrt{u_{k+1}}}\right).
			\end{align*}
			 Suy ra
			\begin{align*}S&=\dfrac{1}{d}\cdot \left(\dfrac{1}{\sqrt{u_1}}-\dfrac{1}{\sqrt{u_2}}\right)+\dfrac{1}{d}\cdot \left(\dfrac{1}{\sqrt{u_2}}-\dfrac{1}{\sqrt{u_3}}\right)+\ldots +\dfrac{1}{d}\cdot \left(\dfrac{1}{\sqrt{u_{2\,023}}}-\dfrac{1}{\sqrt{u_{2\,024}}}\right)\\
				&=\dfrac{1}{d}\cdot \left(\dfrac{1}{\sqrt{u_1}}-\dfrac{1}{\sqrt{u_{2\,024}}}\right)\\ 
				 &=\dfrac{1}{3}\left(1-\dfrac{1}{\sqrt{6\,070}}\right).
			\end{align*}
	\end{itemchoice}}
\end{ex}


%%%%%%%%%%%%%%%% TRẢ LỜI NGẮN %%%%%%%%%%%%%%%%

% nguồn: Ok_TraLoiNgan_Lop11_Lan1
\begin{ex}%[1D2H2-2]%[ma:237ba91d]
	Gọi số  hạng đầu $u_1$ và công sai $d$ của cấp số cộng $\left(u_n \right)$ biết rằng: $u_1+2u_5=0$ và $S_4=14$.
	 Tính $u_1\cdot d$.
	\shortans{-63} 
	 \loigiai{
		Ta có $u_1+2u_5=0\Leftrightarrow u_1+2\left(u_1+4d\right)=0\Leftrightarrow  3u_1+8d=0$. \hfill{(1)}
		\[S_4=14\Leftrightarrow \dfrac{4\left(2u_1+3d\right)}{2}=14\Leftrightarrow 2u_1+3d=7.\; (2)\; 
		\] 
		Giải hệ gồm hai phương trình (1)  và (2), ta được $\heva{&u_1=8\\
		 &d=-3.}$\\
		 Vậy cấp số cộng có số hạng đầu $u_1=8$, công sai $d=-3$.
	 }
\end{ex}

% nguồn: Ok_TraLoiNgan_Lop11_Lan1
\begin{ex}%[1D2H2-2]%[ma:9a3807d3]
	Cho cấp số cộng $(u_n)$ biết $\heva{& u_1-u_2=4 \\& 2u_5+d=-6}$. Tìm số hạng thứ ba của cấp số cộng đã  cho.
	\shortans{7} 
	\loigiai{ 
		Ta có $\heva{& u_1-u_2=4 \\& 2u_5+d=-6}\Leftrightarrow\heva{& u_1-(u_1+d)=4 \\& 2(u_1+4d)+d=-6}\Leftrightarrow\heva{& -d=4 \\& 2u_1+9d=-6}\Leftrightarrow\heva{& d=-4 \\&  u_1=15.}$\\
		Vậy cấp số cộng đã  cho có số hạng đầu $u_1=15$ và công sai $d=-3$.\\
		 Khi đó $u_3=u_1+2d=7$.
	} 
\end{ex}

% nguồn: Ok_TraLoiNgan_Lop11_Lan1
\begin{ex}%[1D2V2-2]%[ma:478671a6]
	Gọi số hạng đầu $u_1$ và công sai $d$ của cấp số cộng $\left(u_n \right)$ biết rằng $\left\{\begin{array}{l}{u_4=10} \\ {u_4+u_6=26}
\end{array}\right.$
Tính  $u_1\cdot d$.
\shortans{3} 
 \loigiai{
	Ta có	\[\begin{array}{ll}\heva{&u_4=10 \\
	&u_4+u_6=26 
	} & \Leftrightarrow  \heva{&u_1+3d=10\\
	&u_1+3d+u_1+5d=26} \\  & \Leftrightarrow \heva{&u_1+3d=10 \\
	&2u_1+8d=26}\Leftrightarrow  \heva{&u_1=1 \\
	&d=3.} 
	\end{array}\]
	Vậy cấp số cộng có số hạng đầu  $u_1=1$ công sai $d=3$.
} 
\end{ex}

% nguồn: Ok_TraLoiNgan_Lop11_Lan1
\begin{ex}%[1D2V2-2]%[ma:3eb70609]
	Giải phương trình sau  $2+7+12+\cdots+x=245$.
	 \shortans{47}
	 \loigiai{
		Ta có dãy số $2,7,12,\ldots, x$ lập thành cấp số cộng  có $\heva{&{u_1=2} \\
		 &{d=5} \\
		&{u_n=x}  \\
		 &{S_n=245}.
		}$\\ 
		 Suy ra $S_n=245\Leftrightarrow 245=\dfrac{n}{2} \left[2u_1+(n-1)d\right]\Leftrightarrow 245.2=n[2.2+(n-1)5]\Rightarrow n=10$.\\
		Vậy  $x=u_{10}=u_1+9d=47$.
	} 
\end{ex}


%%%%%%%%%%%%%%%% TỰ LUẬN %%%%%%%%%%%%%%%%

% nguồn: Ok_TuanLuan_Toan11_Lan2
\begin{ex}%[1D2N2-2]%[ma:ad26d539]
	Cho cấp số cộng $\left(u_{n}\right)$ có số hạng  tổng quát là $u_{n}=3 n-2$. Tìm công sai $d$ của cấp số cộng.
	\loigiai{Ta có $u_{n+1}-u_{n}=3(n+1)-2-3  n+2=3$.\\
	Suy  ra $d=3$ là công sai của cấp số cộng.}
\end{ex}

% nguồn: Ok_TuanLuan_Toan11_Lan2
\begin{ex}%[1D2N2-2]%[ma:36fb1f65]
	Chứng minh các dãy số  sau là cấp số cộng.
	\begin{enumerate}
		\item Dãy số $(u_n)$ với $u_n = 2\,020n - 2\,021$. 
		\item Dãy số $(u_n)$ với $u_n = -2n  + 5$.
		\item Dãy số $(u_n)$ với $u_n  = 3-2n$.
		\item Dãy số $(u_n)$ với $u_n  = 3^n$.
	\end{enumerate}
	 \loigiai{
		\begin{enumerate}
			\item Ta có $u_{n+1} - u_n = 2\,020(n+1) - 2\,021  - (2\,020n - 2\,021) = 2\,020$.\\
			Vậy $(u_n)$ là một cấp số cộng với công sai  $d = 2\,020$.
			\item Ta có $u_{n+1} - u_n = -2(n+1) + 5 - (-2n + 5) =  -2$.\\
			Vậy $(u_n)$ là một cấp số cộng với công  sai $d = -2$.
			\item Ta có $u_1 =  3 - 2.1 = 1;\, u_n - u_{n-1} = 3 - 2n - [3 - 2(n-1)] = -2$ với mọi $n \ge 2$.\\
			Vậy dãy số $(u_n)$ đã cho là một cấp số cộng có số  hạng đầu $u_1 = 1$ và công sai $d = -2$.
			\item Ta  có $u_1 = 3;\, u_2 = 9;\, u_3 = 27;\ldots$ . Do đó, $u_2 - u_1 = 6;\, u_3 - u_2 = 18$.\\
			Vậy dãy số $(u_n)$ đã cho không  là một cấp số cộng.
		\end{enumerate}
	 }
\end{ex}

% nguồn: Ok_TuanLuan_Toan11_Lan2
\begin{ex}%[1D2N2-2]%[ma:03a73b37]
	Chứng minh các dãy số $(u_n)$ sau, dãy số nào không phải là cấp số  cộng.
	\begin{enumerate}
		\item Dãy  số $(u_n)$ với $u_n=n^2+n+1$.
		 \item Dãy số $(u_n)$ với $u_n=(-1)^n+3n$.
	\end{enumerate}
	 \loigiai{
		\begin{enumerate}
			 \item Với $u_n=n^2+n+1$, ta có
			\[ 
			 u_{n+1}-u_n=(n+1)^2+(n+1)+1-(n^2+n+1)=2n+2.
			 \]
			Hiệu số $u_{n+1}-u_n$ phụ thuộc  vào $n$, nên dãy không phải là cấp số cộng.
			\item  Với $u_n=(-1)^n+3n$, ta có
			\[ 
			\begin{aligned}
				u_{n+1}-u_n &=(-1)^{n+1}+3(n+1)-\left[(-1)^n+3n\right]  \\
				 &=-(-1)^n+3-(-1)^n=3-2(-1)^n.
			\end{aligned}
			\] 
			Hiệu số $u_{n+1}-u_n$ phụ thuộc vào $n$,  nên dãy không phải là cấp số cộng.
		\end{enumerate}
	} 
\end{ex}

% nguồn: Ok_TuanLuan_Toan11_Lan2
\begin{ex}%[1D2H2-2]%[ma:647f0549]
	Cho cấp số cộng $\left(u_{n}\right)$ có $u_{1}=-3$, $u_{6}=27$. Tính công sai $d$. 
	\loigiai{Ta có $u_{6}=u_{1}+5  d=27 \Rightarrow d=6$.}
\end{ex}

% nguồn: Ok_TuanLuan_Toan11_Lan2
\begin{ex}%[1D2H2-2]%[ma:ab37947d]
	Cho một  cấp số cộng $\left(u_{n}\right)$ có $u_{1}=\dfrac{1}{3}$, $u_{8}=26$. Tìm công sai $d$.
	 \loigiai{Ta có
		$u_{8}=u_{1}+7 d \Leftrightarrow 26=\dfrac{1}{3}+7  d \Leftrightarrow d=\dfrac{11}{3}$.
	} 
\end{ex}

% nguồn: Ok_TuanLuan_Toan11_Lan2
\begin{ex}%[1D2V2-2]%[ma:1c9c2599]
	Cho cấp số cộng $(u_n)$ với ${u_1=\dfrac{1}{3}}$ và ${u_1+u_2+u_3=-1}$. Tìm công sai $d$ và viết công  thức của số hạng tổng quát $u_n$.
	\loigiai{ 
		Ta có ${u_1+u_2+u_3=-1}$  hay ${3 u_1+3d=-1}$ mà ${u_1=\dfrac{1}{3}}$ nên ${d=-\dfrac{2}{3}}$. \\
		 Công thức của số hạng tổng quát $u_n$ là ${u_n=\dfrac{1}{3}+(n-1)\cdot \left(-\dfrac{2}{3}\right)=-\dfrac{2}{3}n+1}$.
	} 
\end{ex}
